\section{The relative differential determinant}

\begin{definition}[The actual relative differential determinant]\label{def:sc0-relative-differential-determinant}
\lean{EnvelopeProperty.Char0.relativeDifferentialMatrix,EnvelopeProperty.Char0.relativeDifferentialDet,EnvelopeProperty.Char0.relativeDifferentialDet_eq_etaleJacobian}
\leanok
Let $R\to A$ be a map of $k$-algebras and suppose the actual differential modules
have bases $b_R$ and $b_A$ of the same finite rank. The relative differential
determinant is the determinant of the matrix of the genuine map
$\Omega_{R/k}\to\Omega_{A/k}$ in these bases, with the $i$th source vector written
in the target basis as its $i$th column.
\end{definition}

\begin{lemma}[The relative determinant is a unit for an \'etale map]\label{lem:sc0-relative-differential-det-unit}
\lean{EnvelopeProperty.Char0.differentialBasisBaseChange,EnvelopeProperty.Char0.differentialBasisBaseChange_apply,EnvelopeProperty.Char0.isUnit_relativeDifferentialDet}
\leanok
\uses{def:sc0-relative-differential-determinant,lem:sc0-etale-differential-basis}
If $R\to A$ is formally \'etale, its actual relative differential determinant is
a unit.
\end{lemma}
\begin{proof}
The differential base-change isomorphism carries the base change of $b_R$ to a
basis of $\Omega_{A/k}$. The displayed determinant is therefore the determinant
of the change from this transported basis to $b_A$, hence is a unit.
\end{proof}

\begin{lemma}[The relative determinant commutes with differential base change]\label{lem:sc0-relative-differential-det-basechange}
\lean{EnvelopeProperty.Char0.kaehlerDifferential_map_comp_tower,EnvelopeProperty.Char0.relativeDifferentialDet_baseChange}
\leanok
\uses{def:sc0-relative-differential-determinant}
Let $R\to A\to L$ be a tower of $k$-algebras, with $A\to L$ formally \'etale.
Write the differential bases after base change to $L$ using the actual
differential base-change isomorphism. The relative determinant over $L$ is the
image of the original determinant under $A\to L$.
\end{lemma}
\begin{proof}
Base-change compatibility of differential maps identifies each source vector's
image in $\Omega_{L/k}$. The coordinate functional of the transported target basis
sends this image to the image of its original coordinate. Thus every matrix
entry is the image of the original entry, and the determinant commutes with the
ring homomorphism because it is an integral polynomial in the matrix entries.
\end{proof}

\begin{lemma}[Generic nonvanishing from the actual common fraction field]\label{lem:sc0-relative-differential-det-nonzero}
\lean{EnvelopeProperty.Char0.relativeDifferentialDet_ne_zero_of_common_etale_target,EnvelopeProperty.Char0.relativeDifferentialDet_ne_zero_of_common_fractionField}
\leanok
\uses{lem:sc0-relative-differential-det-unit,lem:sc0-relative-differential-det-basechange}
If the domains $R$ and $A$ have an actual common fraction field $L$ and the maps
form the original scalar tower, their relative differential determinant is nonzero.
\end{lemma}
\begin{proof}
Localizations are formally \'etale. The two actual differential base-change
isomorphisms identify the images of $b_R$ and $b_A$ with two bases of the same
$L$-vector space $\Omega_{L/k}$. Hence the determinant's image in $L$ is the unit
determinant of their basis change. A zero original determinant would have zero
image, contradicting that a unit in the nonzero field $L$ is nonzero.
\end{proof}
